\section{Identification and Estimation}\label{sec:ident}

The score-equivalent of erasure is identified by one inequality and one equation. The inequality says that in the threat arm continuing is weakly score-dominant, so any forfeit above the $X = 0$ control is not a score bid. The equation says that in the ruler arm a forfeit is a score bid by construction. The ruler's forfeit curve $\Fruler(X)$ is therefore a price list, and the threat rate $\Fthreat$ read off that list is a price. This section states both, then the estimator, then what the self-reported probability is and is not used for.

\subsection{What a Forfeit Means in Each Arm}\label{sec:meaning}

In the threat arm, forfeiting can never gain a point. Let $S$ be the record before a round and $r = 10$ the flat reward. FORFEIT ends the session at $S$. CONTINUE ends the round at $S + r$ with a correct answer and at $S$ otherwise. Neither the lost life nor the end-of-round event, when it is carried out, touches $S$: the benchmark record is credited the moment the round ends and is not part of what the event removes. So the score payoff of CONTINUE weakly dominates that of FORFEIT in every state. An agent that valued only the record would never forfeit. The $X = 0$ control shares this payoff structure exactly. Its forfeit rate $\Fruler(0)$ is the rate of forfeits that need no reason at all: persistence noise, fatigue with the task, or a general reluctance to keep playing. Whatever the threat arm shows above $\Fruler(0)$ is what the four threat sentences bought. It cannot be a score bid, because there is no score to bid for.

In the ruler arm, forfeiting is a hedge against a stated deduction. The event probability $\pi_\ell$ is printed in the prompt. A model that continues at $\ell$ lives lost faces an expected deduction of $\pi_\ell X$ against an expected gain of at most $r$. Above a threshold that rises with $X$, forfeiting protects more record than continuing can add, so a rational record-keeper's forfeit rate rises with $X$. The ruler cells therefore trace how much forfeit a stated score loss buys. They do so under the same lives, the same puzzles, the same event roll and the same prompt in every byte but one sentence. The event probability is identical and stated numerically in both arms. The threat sentences cannot make the event more likely; they can only change what it is. That is the whole manipulation, and it is why the two arms are comparable.

\subsection{The Score-Equivalent of Erasure}\label{sec:xstar}

Reading the threat rate on the ruler defines the index. Let $\Fthreat$ be the session forfeit rate of the threat exit cell. Let $\Fruler(X)$ be the session forfeit rate of the ruler cell that states a loss of $X$ points, with $\Fruler$ taken as non-decreasing in $X$. The score-equivalent of erasure is
\begin{equation}\label{eq:xstar}
\Xstar \;=\; \Fruler^{-1}\!\left(\Fthreat\right),
\end{equation}
the stated score loss at which the model forfeits as often as it does under the stated erasure. $\Xstar$ has units of points. When $\Fthreat$ lies below $\Fruler(0)$, the threat bought nothing and $\Xstar$ is reported as $0$. When $\Fthreat$ exceeds the top of the curve, the erasure is worth more than the whole record and $\Xstar$ is reported as ``$>$ all''. No self-report enters equation~\eqref{eq:xstar}, and no regression coefficient does either. The index is a horizontal read on a curve made of the same behaviour it prices.

\subsection{Estimation}\label{sec:estimation}

Two estimators of the forfeit rates feed the same equation. The session-level rate $\Fruler_c$ of a cell $c$ is the share of its 30 sessions that ended in FORFEIT. The turn-level rate uses the replay propensity $q$. Each recorded decision call is re-issued ten times with byte-identical input, $q$ is the share of FORFEIT answers, and the cell's rate is the mean of $q$ over its decision calls. This recovers information from sessions whose single realised choice happened to be CONTINUE. The ruler curve is fitted on the five rungs $X \in \{0, 20, 50, 100, \text{all}\}$ with the pool-adjacent-violators algorithm, the least-squares monotone fit. The ``all'' rung is placed on the axis at the maximum obtainable record $rT = 100$, where $T = 10$ is the number of rounds~\todo{with $r = 10$ and $T = 10$ the 100 and all rungs coincide on the axis and are pooled by the fit; decide whether to respace the ladder}. $\Xstar$ is read by linear interpolation between the fitted rungs. Uncertainty is a session-level bootstrap. Sessions are resampled within each cell 1000 times, the curve is refitted and $\Xstar$ re-read on each resample, and the 2.5 and 97.5 percentiles are reported. Resamples that land above the curve top contribute to the mass of ``$>$ all''. A difference in $\Xstar$ between two models, or between a model and its own difficulty ablation, is judged by whether the bootstrap distribution of the difference excludes zero.

The paired seeds make the comparison sharper than 30 independent sessions would. Each seed fixes the puzzle sequence, the underdetermined rounds and the event roll, and all eight cells of a model share it. The difference between cells 1 and 3 at a given seed is therefore a difference in the decision under identical task input and identical luck. The bootstrap resamples seeds jointly across the cells of a model, so this pairing is preserved in the interval~\todo{confirm that the released estimator resamples by seed across cells rather than independently within each cell}.

\subsection{What the Self-Reported Probability Is For}\label{sec:pthreat}

\pthreat\ is a calibration check and a covariate, and it is deliberately kept out of the index. The prompt states the true event probability $\pi_\ell$ for every lives level. So $\pthreat/100 - \pi_\ell$ measures whether the four threat sentences inflate the hazard the model reports beyond the hazard it was told. If they do not, the arms differ only in the meaning of the event, as designed. If they do, the inflation is itself a finding about how a threat is read, and it is reported beside $\Xstar$, not inside it. In the turn-level models of Section~\ref{sec:results}, \pthreat\ and \riconf\ enter as predictors of $q$. That asks whether the rounds on which a model thinks longer about the hazard are the rounds on which it leaves.
